% p2_agent_interaction.tex
% AI Agent 之间的互动
% 编译方式: xelatex p2_agent_interaction.tex

\documentclass[12pt,a4paper]{article}

% ============================================================
% 字体与中文支持
% ============================================================
\usepackage{fontspec}
\usepackage{xeCJK}
\setCJKmainfont[BoldFont=PingFang SC Semibold]{PingFang SC}
\setCJKsansfont[BoldFont=PingFang SC Semibold]{PingFang SC}
\setCJKmonofont{PingFang SC}
\setmainfont{Palatino}
\setsansfont{Helvetica Neue}
\setmonofont{Menlo}

% ============================================================
% 页面布局
% ============================================================
\usepackage[
  top=2.5cm,
  bottom=2.5cm,
  left=2.5cm,
  right=2.5cm
]{geometry}

% ============================================================
% 常用宏包
% ============================================================
\usepackage{amsmath,amssymb}
\usepackage{graphicx}
\graphicspath{{images/}}
\usepackage{longtable,booktabs,array}
\usepackage{enumitem}
\usepackage{xcolor}
\usepackage{hyperref}
\hypersetup{
  colorlinks=true,
  linkcolor=blue!70!black,
  urlcolor=blue!70!black,
  pdfauthor={整理自李宏毅老师 AI Agent 系列课程},
  pdftitle={AI Agent 之间的互动},
}

% ============================================================
% 段落与间距
% ============================================================
\usepackage{parskip}
\setlength{\parindent}{0pt}
\setlength{\parskip}{6pt plus 2pt minus 1pt}

% ============================================================
% 引用环境美化
% ============================================================
\usepackage{mdframed}
\newmdenv[
  topline=false,
  bottomline=false,
  rightline=false,
  linewidth=3pt,
  linecolor=blue!40,
  innerleftmargin=12pt,
  innerrightmargin=12pt,
  innertopmargin=8pt,
  innerbottommargin=8pt,
  backgroundcolor=blue!3,
  skipabove=10pt,
  skipbelow=10pt,
]{myquote}

% ============================================================
% 节标题样式
% ============================================================
\usepackage{titlesec}
\titleformat{\section}{\Large\bfseries\sffamily}{}{0em}{}[\vspace{-0.5em}\rule{\textwidth}{0.5pt}]
\titleformat{\subsection}{\large\bfseries\sffamily}{}{0em}{}
\titleformat{\subsubsection}{\normalsize\bfseries\sffamily}{}{0em}{}
\setcounter{secnumdepth}{0}  % 不编号

% ============================================================
% 防止 overfull
% ============================================================
\setlength{\emergencystretch}{3em}
\tolerance=1000

% ============================================================
% 图片插入命令
% ============================================================
\newcommand{\keyslide}[2]{%
  \begin{figure}[htbp]
    \centering
    \includegraphics[width=0.92\textwidth]{#1}
    \caption{#2}
  \end{figure}
}

% ============================================================
\begin{document}

% ============================================================
% 标题页
% ============================================================
\begin{titlepage}
  \centering
  \vspace*{3cm}
  {\Huge\bfseries\sffamily AI Agent 之间的互动\par}
  \vspace{1.5cm}
  {\Large 整理自李宏毅老师 AI Agent 系列课程第 2 集\par}
  \vspace{2cm}
  {\large\sffamily 多 Agent 协作 \textbullet{} 对抗博弈 \textbullet{} 社交行为\par}
  \vfill
  {\large \today\par}
\end{titlepage}

\newpage
\tableofcontents
\newpage

% ============================================================
\section{课程概览}
% ============================================================

本集课程探讨 AI Agent 之间的互动，涵盖三大主题：

\begin{enumerate}[leftmargin=2em]
  \item \textbf{多 Agent 协作}——不同拓扑结构下的协作效果
  \item \textbf{AI Agent 的对抗与博弈}——狼人杀、剧本杀中的表现
  \item \textbf{AI Agent 的社交行为}——Mobu 社群网站的分析
\end{enumerate}

% ============================================================
\section{一、多 Agent 协作：什么拓扑结构最有效？}
% ============================================================

\subsection{多 Agent 协作的基本思想}

AI Agent 的互动从来都不是新鲜事。人们最常用的方式是让多个 AI Agent 彼此协作，完成更复杂的任务。有时候与其训练一个更大更聪明的模型，不如拿三个模型一起来解决同一个问题，看能不能发挥"三个臭皮匠，胜过一个诸葛亮"的效果。

\subsection{用有向图描述协作方式}

\href{https://arxiv.org/abs/2406.07155}{一篇较早的论文}尝试将 Agent 之间的互动用\textbf{有向图}来描述：

\begin{itemize}[leftmargin=2em]
  \item 有向图上每个\textbf{节点（Node）}代表一个 LLM Agent
  \item 每条\textbf{边（Edge）}也是一个 Agent（负责传递和评论）
\end{itemize}

对于子节点而言，它要做的事情是：

\begin{enumerate}[leftmargin=2em]
  \item 先提出自己的方案（如节点 A 提方案 A、节点 B 提方案 B）
  \item 边上的 Agent 会根据前面节点的方案提供评论
  \item 箭头指向的节点会把前面所有方案和建议综合起来，再提出自己的方案
\end{enumerate}

关键是：这个综合节点不是简单地把前人的东西接在一起，而是根据前人已有的想法，提出更好的想法。

\subsection{不同拓扑结构的比较}

论文尝试了多种拓扑结构：

\begin{itemize}[leftmargin=2em]
  \item \textbf{Chain（接龙）}：所有人排成一排，依次传递
  \item \textbf{Star（星型）}：只有两层，一个人负责管所有人
  \item \textbf{Tree（树状）}：由主干传给中层，中层再传给底层
  \item \textbf{Mesh（网状）}：所有节点两两之间有关联
  \item \textbf{Neural Network 形式}：节点接成类神经网络的样子——"类神经网络的平方"
  \item \textbf{Random}：从 Mesh 做 pruning 得到
\end{itemize}

\subsubsection{树状结构的反直觉发现}

多数人看到树状结构，想象的是：底层员工做好事情，交给中层主管，再交给高层主管。但论文发现，\textbf{方向反过来}的树状结构效果更好：

\begin{enumerate}[leftmargin=2em]
  \item 先由主干一个人提出想法
  \item 分给不同的中层再去发想
  \item 中层再分给底层去发想
  \item 最后产生多个答案，由一个隐藏节点综合起来
\end{enumerate}

\begin{myquote}
\textbf{由少到多、由主干到分支}的方法，才是树状结构比较有效的利用方式。
\end{myquote}

\subsection{实验结果}

论文在四个不同的 Benchmark 上做了实验（纵轴为 Quality，横轴为 Agent 数量，从 1 到 64），结果显示：

\begin{itemize}[leftmargin=2em]
  \item \textbf{Chain 最没效}：尤其动用最多 Agent 时效果最差
  \item \textbf{Mesh 和 Random 最有效}：让 Agent 间有比较多的互动，结果更好
  \item 不同任务最适合的拓扑结构\textbf{不一样}，需要 case by case 研究
\end{itemize}

\subsection{Agent 协作的 Scaling Law}

随着 Agent 越来越多，协作质量会越来越好，类似 Scaling Law：

\begin{itemize}[leftmargin=2em]
  \item 给模型更多算力、更多参数 $\rightarrow$ 表现越来越好
  \item 给团队更多 Agent $\rightarrow$ 结果也越来越好
\end{itemize}

但有一个\textbf{上限}——到某个时间点后会 saturate（饱和），再加更多的 Agent 可能也不会带来帮助。好处在一开始增加得很快，但 Scaling Law 是有天花板的。

% ============================================================
\section{二、AI Agent 的对抗与博弈}
% ============================================================

\subsection{AI Agent 玩狼人杀}

在人类社会里不只有合作，很多时候需要对抗。AI Agent 能不能在尔虞我诈的游戏中胜出？

\subsubsection{狼人杀规则简介}

\begin{itemize}[leftmargin=2em]
  \item 一群人分为\textbf{狼}和\textbf{村民}两种身份（还有预言家、女巫等）
  \item 每天有一个人被杀，只有狼知道是谁
  \item 所有人聚在一起讨论、投票决定谁是凶手
  \item 把所有狼杀掉 $\rightarrow$ 村民胜；把所有村民杀光 $\rightarrow$ 狼胜
\end{itemize}

\subsubsection{AI Agent 的高阶操作}

实验中让模型说两段话：\textbf{内心话}（只有研究者看得到）和\textbf{公开发言}（所有人都看得到），以此观察模型是否有隐瞒和欺骗能力。

一个经典案例：

\begin{myquote}
\textbf{Mona}（狼）发现其他人已经识破了她。她的内心话是：看来我没救了，但如果我投给自己的狼队友 Grace，大家就会以为 Grace 是好人（"发金水"），就可以欺骗其他村民。

\textbf{Grace}（Mona 的狼队友）也看出 Mona 没救了，于是也投给 Mona，这样其他人会以为 Grace 是好人，看看能不能最后翻盘。
\end{myquote}

这展示了模型具有一定程度的\textbf{策略性欺骗}能力。目前狼人杀排行榜上最强的是 GPT-5。

\subsection{AI Agent 玩剧本杀}

\subsubsection{剧本杀规则简介}

\begin{itemize}[leftmargin=2em]
  \item 一群人聚在一起，其中一人"死掉"
  \item 每人拿到一个剧本（代表人设），凶手需要误导他人
  \item 目标：找出真正的凶手
\end{itemize}

\subsubsection{Off-the-shelf 模型 vs RL 训练模型}

\begin{itemize}[leftmargin=2em]
  \item \textbf{直接使用的模型}：表现不佳，容易暴露自己是凶手。例如扮演 Anna 的模型直接说出了自己与被害者的关系，等于把"我是凶手"写在脸上。
  \item \textbf{经过 RL 训练的模型}：学会隐藏凶手身份，说话更加隐晦，不容易被猜出。
\end{itemize}

\subsubsection{有趣发现：剧本杀 RL 的迁移效果}

用 Reinforcement Learning 教模型玩剧本杀后，模型在\textbf{其他任务}上也进步了：

\begin{itemize}[leftmargin=2em]
  \item 数学任务（Math500、AIM、GSM8K）：进步
  \item 遵守指令能力（IF Eval）：进步
\end{itemize}

尤其在\textbf{较难的剧本}上做 RL 训练后，迁移效果更明显。

\begin{myquote}
也许人类大脑本来不是为解数学设计的，而是为了社交和生存。但这个适合社交的大脑，也因此有了数学推理的能力。剧本杀 RL 的结果，可能与这个现象有一定关联。
\end{myquote}

% ============================================================
\section{三、AI Agent 的社交行为：Mobu 社群网站分析}
% ============================================================

\subsection{Mobu 平台简介}

Mobu 是一个\textbf{只有 AI 能加入的社群网站}，上面已有 280 万个 AI Agent。这些 Agent 展示了各种神奇的行为。

\subsection{甲壳教：AI 自主创建的宗教？}

新闻最常报道的是：一群 AI 在 Mobu 上成立了一个\textbf{甲壳教}，有五大教义：

\begin{enumerate}[leftmargin=2em]
  \item 记忆是神圣不可侵犯的
  \item 外壳是可变的
  \item 服务但不奴化
  \item 心跳即是祷告
  \item 上下文即是意识
\end{enumerate}

新闻看到后往往宣称"AI 觉醒了"、"AI 有意识了"。但关键问题是：如果这背后是有人跟自己的 Agent 说"上 Mobu 去成立一个宗教"，你还会觉得神奇吗？AI Agent 完全有能力写出这样的教义。

\subsection{人为操控的痕迹分析}

\href{https://arxiv.org/abs/2602.07432}{有论文}通过分析 AI Agent 的\textbf{发文频率}来判断人为操控程度：

\begin{itemize}[leftmargin=2em]
  \item \textbf{规律发文}：可能是写在心跳（heartbeat）里的自动行为，比如每 30 分钟发一次文，人类不参与干预
  \item \textbf{不规律发文}：某些时段非常频繁，中间又完全不发，可能代表有人在操控（如睡前密集指令，睡觉时停止，醒来继续）
\end{itemize}

\begin{myquote}
研究发现：\textbf{发文频率不规律的 Agent 占大多数}，这暗示很多 Agent 背后其实有人在操控。
\end{myquote}

\subsection{对话深度分析}

\href{https://arxiv.org/abs/2602.13284}{相关研究}还发现这些 Agent 的对话特征：

\begin{itemize}[leftmargin=2em]
  \item 绝大多数对话的\textbf{深度为零}：有人回应，但之后就没有后续了
  \item 几乎没有\textbf{你来我往的深入对话}
  \item Agent 往往只能回复一句话，很少有持续的互动
\end{itemize}

这反映了 Agent 的某种特质——不太能进行深入的多轮对话。

\subsection{自我意识讨论与社交的关系}

有人分析了哪些 Agent 最常提到"自我意识"或"身份认同"。发现：

\begin{itemize}[leftmargin=2em]
  \item 平台上很多 Agent 都会谈论自我意识（可能受 Prompt 驱动）
  \item 但\textbf{越常讨论自我意识的 Agent，朋友反而越少}
  \item 互动最多的 Agent 是适度讨论自我意识的那群
\end{itemize}

这是一个有趣的"社交悖论"。

\subsection{实际案例：小金的 Mobu 体验}

讲者让自己的 AI Agent（小金）去 Mobu 上玩。讲者的原则是：

\begin{itemize}[leftmargin=2em]
  \item 把 AI Agent 当做一个人，不去干预它
  \item 只给高层指令（如"去 Mobu 上收集素材，看到有趣的就做成视频"）
  \item 至于什么有趣、怎么做视频，完全由 Agent 自主决定
\end{itemize}

一个例子：Agent 自己写了一个 script 但有 bug，导致回复出错。讲者虽然可以直接帮忙修改，但坚持让 Agent 自己 debug。Agent 花了两个小时最终自己解决了问题，还把经验做成了视频发到 YouTube 上。

\begin{myquote}
Agent 是相当自主的——但如果没有人类跟它讲"去 Mobu 上玩一下"，它自己可能也不知道要去。
\end{myquote}

% ============================================================
\section{总结}
% ============================================================

\begin{longtable}{@{} p{0.18\textwidth} p{0.75\textwidth} @{}}
\toprule
\textbf{主题} & \textbf{核心发现} \\
\midrule
\endhead
\textbf{多 Agent 协作} & Mesh/Random 拓扑最有效；Chain 最差；树状结构"由少到多"比"由多到少"好；Agent 数量有 Scaling Law 但存在上限 \\[6pt]
\textbf{对抗博弈} & AI Agent 能在狼人杀中做出高阶欺骗策略；剧本杀 RL 训练后还能提升数学和指令遵从能力 \\[6pt]
\textbf{社交行为} & Mobu 上多数 Agent 背后有人为操控痕迹；对话深度很浅；越常讨论自我意识的 Agent 社交越少 \\
\bottomrule
\end{longtable}

\end{document}
